\documentclass[11pt]{article}
\usepackage[margin=1in]{geometry}
\usepackage{booktabs}
\usepackage{graphicx}
\usepackage{amsmath,amssymb}
\usepackage[round,authoryear]{natbib}
\usepackage{hyperref}
\usepackage{microtype}
\usepackage{enumitem}
\usepackage{xcolor}
\usepackage{listings}
\usepackage{fontspec}
\newfontfamily\emojifont{Apple Color Emoji}[Renderer=Harfbuzz]
\newcommand{\parrotemoji}{{\emojifont 🦜}}
\newcommand{\given}{\,\vert\,}
\newcommand{\twin}{\widehat P}
\newcommand{\zwill}{\texttt{zwill}}
\DeclareMathOperator*{\E}{\mathbb E}
\DeclareMathOperator{\KL}{KL}
\definecolor{codegray}{RGB}{246,247,249}
\definecolor{codegreen}{RGB}{25,105,60}
\definecolor{codeblue}{RGB}{30,76,145}
\definecolor{codeorange}{RGB}{166,74,0}
\lstdefinestyle{shell}{
  language=bash,
  basicstyle=\ttfamily\footnotesize,
  backgroundcolor=\color{codegray},
  frame=single,
  rulecolor=\color{black!18},
  breaklines=true,
  columns=fullflexible,
  keepspaces=true,
  showstringspaces=false,
  commentstyle=\color{codegreen},
  keywordstyle=\color{codeblue}\bfseries,
  stringstyle=\color{codeorange},
  xleftmargin=0.5em,
  xrightmargin=0.5em,
  aboveskip=0.7em,
  belowskip=0.7em
}
\lstdefinestyle{prompt}{
  basicstyle=\ttfamily\scriptsize,
  backgroundcolor=\color{codegray},
  frame=single,
  rulecolor=\color{black!18},
  breaklines=true,
  columns=fullflexible,
  keepspaces=true,
  showstringspaces=false,
  xleftmargin=0.5em,
  xrightmargin=0.5em
}
\lstdefinestyle{json}{
  basicstyle=\ttfamily\scriptsize,
  backgroundcolor=\color{codegray},
  frame=single,
  rulecolor=\color{black!18},
  breaklines=true,
  breakatwhitespace=true,
  columns=fullflexible,
  keepspaces=true,
  showstringspaces=false,
  upquote=true,
  xleftmargin=0.5em,
  xrightmargin=0.5em,
  aboveskip=0.7em,
  belowskip=0.7em
}
\newcommand{\TotalPredictions}{36,600}
\newcommand{\TotalCost}{\$523.30}
\newcommand{\SimNLLPass}{14}
\newcommand{\SimNLLFail}{0}
\newcommand{\SimNLLInconclusive}{0}
\newcommand{\ByQNLLPass}{13}
\newcommand{\ByQNLLFail}{0}
\newcommand{\ByQNLLInconclusive}{1}
\newcommand{\UniformQuestionsBetter}{65}
\newcommand{\UniformNLLDelta}{-0.312}
\newcommand{\UniformNLLLo}{-0.364}
\newcommand{\UniformNLLHi}{-0.259}
\newcommand{\MarginalMAE}{0.048}
\newcommand{\MarginalCorrelation}{0.925}
\newcommand{\MarginalTopAgreement}{79.7\%}
\newcommand{\SemanticNLLCoef}{-0.995}
\newcommand{\SemanticNLLLo}{-1.860}
\newcommand{\SemanticNLLHi}{0.068}
\newcommand{\SemanticAdvCoef}{0.236}
\newcommand{\SemanticAdvLo}{-2.822}
\newcommand{\SemanticAdvHi}{3.243}
\newcommand{\GSSQuestions}{38}
\newcommand{\GSSCalls}{7,638}
\newcommand{\GSSCost}{\$73.97}
\newcommand{\GSSConditionalDelta}{-0.333}
\newcommand{\GSSConditionalLo}{-0.432}
\newcommand{\GSSConditionalHi}{-0.251}
\newcommand{\GSSOneShotDelta}{-0.190}
\newcommand{\GSSOneShotLo}{-0.237}
\newcommand{\GSSOneShotHi}{-0.148}
\newcommand{\GSSMarginalCorrelation}{0.955}
\newcommand{\GSSMarginalMAE}{0.044}
\newcommand{\KALRORespondents}{832}
\newcommand{\KALROQuestions}{20}
\newcommand{\KALROCalls}{4,020}
\newcommand{\KALROCost}{\$39.76}
\newcommand{\KALROConditionalDelta}{-0.351}
\newcommand{\KALROConditionalLo}{-0.609}
\newcommand{\KALROConditionalHi}{-0.127}
\newcommand{\KALROConditionalWins}{13}
\newcommand{\KALROOneShotDelta}{-0.007}
\newcommand{\KALROOneShotLo}{-0.039}
\newcommand{\KALROOneShotHi}{0.027}
\newcommand{\KALROMarginalDelta}{0.102}
\newcommand{\KALROMarginalCorrelation}{0.763}
\newcommand{\KALROMarginalMAE}{0.177}

\newcommand{\OneShotCalls}{69}
\newcommand{\OneShotCost}{\$0.70}
\newcommand{\OneShotMAE}{0.059}
\newcommand{\OneShotTwinMAE}{0.047}
\newcommand{\OneShotNLLDelta}{-0.186}
\newcommand{\OneShotTwinWins}{62}
\newcommand{\OneShotBatteryWins}{14}

\title{\protect\zwill: Reproducible Construction and Validation of Survey Digital Twins}
\author{John J. Horton\\
\normalsize MIT $\cdot$ NBER $\cdot$ E[\parrotemoji]}
\date{\today\\[1em]
\includegraphics[width=0.22\textwidth]{../../../docs/assets/zwill-package.png}}
\begin{document}
\maketitle

\begin{abstract}
Language models make it easy to generate plausible survey respondents but do
not make it easy to establish that those representations predict the right
answers for the right people. Digital-twin evaluations can be distorted by
target leakage, unequal information sets, changes in the scored sample,
provider-specific execution code, silent parsing failures, or benchmarks that
test only population prevalence. We introduce \zwill, an open-source
package that makes the path from raw survey microdata to validated digital-twin
evidence explicit and reproducible. \zwill{} preserves provenance, constructs and
registers holdout experiments, renders auditable respondent prompts, imports
model results, fits exact-information statistical baselines, checks leakage,
computes uncertainty, and builds reports. Expected Parrot provides the portable
multi-provider inference layer: \zwill{} exports serialized EDSL jobs, Expected
Parrot executes them remotely and preserves raw Results and cost records, and
\zwill{} reimports those artifacts for validation. We demonstrate the system on
15 public microdata batteries and then extend the leave-one-question-out design
to five independently sourced 2024 General Social Survey batteries. Across two holdout designs and a three-model
robustness exercise, \TotalPredictions\ respondent-level forecasts all parse
successfully. On 14 clean batteries, GPT-5.5 twins beat matched conditional
XGBoost in all 14 under simultaneous holdout and in 13 under by-question
holdout, with one inconclusive result. A separate same-model one-shot study
shows that respondent context improves NLL on \OneShotTwinWins{} of 69 questions.
In the GSS extension, twins beat exact-context XGBoost on 37 of 38 questions
and the same-model no-persona forecast on 36 of 38.
An external validation on 20 questions from a Kenyan agricultural field study
again favors twins over exact-context XGBoost, while twins and same-model
one-shot forecasts are statistically indistinguishable overall.
These are finite-sample predictive benchmarks, not evidence that synthetic
responses are interchangeable with human observations for downstream causal
or descriptive inference. The empirical results illustrate the package's
purpose: credible claims require an auditable
chain of data, information-set, execution, baseline, and reporting contracts,
not merely convincing synthetic answers.
\end{abstract}

\section{Introduction}
Language models make it inexpensive to construct digital representations of
survey respondents. They do not make it easy to know whether those
representations work. A model can produce persuasive answers by reproducing
population stereotypes, exploit information inadvertently revealed elsewhere
in a questionnaire, or appear superior to a benchmark that received less
information. These failures are difficult to detect from final predictions
alone. Credible validation requires control over the entire path from raw
microdata to prompts, model execution, matched baselines, audits, uncertainty
estimates, and reports.

We introduce \zwill, an open-source command-line package for constructing
and validating survey digital twins. The package treats evaluation as an
experimental-design and evidence-management problem rather than as a prompt
demonstration. It records the respondents and questions under study, the answer
held out from each twin, the exact context permitted for each prediction, the
rendered prompt, the raw model result, and the cases ultimately scored. It then
compares twins with baselines ranging from uniform prediction to a conditional
XGBoost model trained on the same respondents and the exact information exposed
in the corresponding prompts.

This machinery is needed because seemingly small implementation choices can
change the estimand. A statistical model trained on the full questionnaire is
not a fair comparator for a twin shown eight answers. A target-specific
follow-up can reveal a nominally held-out response. Complete-case filtering can
give two methods different respondents. An empirical target marginal is a
useful oracle but not an out-of-sample forecast. Dropping malformed model rows
can make the reported evaluation sample depend on model behavior. \zwill{} turns
these choices into explicit, inspectable artifacts and blocks its configured
benchmark gate when the required matched baseline or prompt-leakage audit is
unavailable. Passing that gate is evidence about the benchmark, not a license
to substitute synthetic observations for humans in arbitrary analyses.

\subsection{\zwill{} and Expected Parrot}
Reliable digital-twin research requires both a scientific workflow and a
reproducible execution layer. \zwill{} owns the former: ingestion, codebooks,
provenance, holdout design, twin construction, approval, leakage policy,
baseline matching, scoring, uncertainty, and reporting. Expected Parrot owns
the latter. \zwill{} serializes the approved experiment as an EDSL job package;
Expected Parrot executes that package remotely across model providers,
preserves the raw Results object, and records model, usage, and cost metadata;
\zwill{} then imports the immutable result package and applies the registered
validation procedure.

The separation is substantive rather than cosmetic. Provider-specific API
loops can alter parameters, error handling, and retained outputs across models.
A portable job/result boundary keeps experimental design out of execution code,
makes multi-provider comparisons use the same scenarios, and leaves a durable
record of what was purchased. In this paper, every paid inference operation is
represented by a \zwill-generated Expected Parrot job package and invoked by a
Makefile target. Every reported table and figure is regenerated from imported
\zwill{} evidence bundles.

The paper makes three contributions. First, it specifies a fair
leave-one-question-out estimand for respondent-level digital twins and a ladder
of baselines that separates general population knowledge from individual
signal. Second, it provides an executable software architecture implementing
the complete validation chain. Third, it demonstrates that architecture across
public microdata batteries, including cases in which leakage audits or strong
benchmarks qualify an apparently favorable result. The empirical application
is therefore both evidence about digital twins and a worked test of the package.

The primary estimand is predictive performance for a previously unseen
question among respondents drawn from the same survey battery and sampling
frame. It concerns individual predictive distributions, not merely recovery of
aggregate marginals. A positive result therefore requires improvement over a
conditional statistical baseline evaluated on the same scored cases. The
experiment does not estimate causal effects, simulate treatment responses, or
validate arbitrary counterfactual behavior.

\subsection{Why leave-one-question-out validation?}
A digital twin is useful only to the extent that it predicts information not
used to construct it. Evaluating answers already supplied in the persona would
measure recall, and comparing only aggregate shares would allow a model to
succeed by reproducing a population stereotype. Prospective validation on
entirely new surveys would be attractive, but then changes in wording,
population, field period, and sampling frame are confounded with twin quality.
Leave-one-question-out validation occupies the useful middle ground: the
answer is genuinely hidden, the realized outcome is known, and respondents,
instrument, and field conditions are held fixed.

The design is the survey analogue of cross-validation. For question $j$, a
construction function $g$ maps permitted respondent information
$C_{ij}=\{x_i,y_{i,-j}\}$ to a probabilistic forecast,
\[
 \widehat p_{ij}=g(q_j,\mathcal A_j,C_{ij}),
 \qquad y_{ij}\notin C_{ij},
\]
which is scored against the answer that was sequestered before construction.
Repeating this operation across questions tests transfer rather than fitting a
question-specific response model. The by-question design most closely matches
the deployment question ``given everything already known about this person,
how would they answer one new item?'' The simultaneous design is a stricter
stress test in which a whole target set is absent. Neither design establishes
external validity to new populations or causal counterfactuals, but both give
an internally valid test of incremental respondent-level signal.

\subsection{Validation target and epistemic scope}
Hullman et al. distinguish heuristic ``validate-then-simulate'' exercises from
statistically calibrated uses of model outputs \citep{hullman2026human}. Their
critique applies directly to how our results should be read. A successful
holdout establishes that a fixed construction predicts observed human answers
on these respondent--item pairs. It does not establish that model-generated
answers may replace human observations in a new study. Calling both activities
``validation'' obscures a change in estimand.

Following \citet{ludwig2025llm}, separate a prediction problem from a
measurement-and-estimation problem. Our prediction estimand is the risk of a
frozen predictor $f$ over the evaluated distribution of respondent--item
inputs $X=(q,C)$:
\[
 R_{\mathcal T}(f)=\E_{(X,Y)\sim\mathcal T}
       [\ell\{Y,f(X)\}],
\]
where $\mathcal T$ is, conservatively, the finite collection of batteries,
targets, and sampled respondents analyzed here. The human answer $Y$ exists for
every scored prediction. The matched XGBoost comparison asks whether one
predictor has lower risk than another on exactly those cases.

The holdout $Y\notin C$ prevents direct target leakage from the constructed
persona, but it is not Ludwig et al.'s ``no training leakage'' condition. These
are public survey items, and proprietary providers do not reveal training
corpora. Question wording, reports, or response distributions may have entered
pretraining. Consequently, we do not estimate performance over an unseen
population of scenarios. A defensible scenario-level generalization would
require defining that population and randomly sampling held-out scenarios, or
using a model whose training boundary makes overlap testable. Our bootstrap
intervals quantify respondent sampling uncertainty conditional on the selected
questions and frozen pipeline; they do not solve training leakage or turn the
69 targets into a random sample of all possible survey questions.

Downstream substitution requires a different condition. Suppose a researcher
uses twin output $\widehat Y$ to estimate a human parameter $\theta_0$ defined
by
\[
 \E[g(Y,X;\theta_0)]=0.
\]
Low NLL, high accuracy, and good aggregate marginal fit do not imply
\[
 \E[g(\widehat Y,X;\theta)]
   =\E[g(Y,X;\theta)]\quad\text{for all }\theta.
\]
For example, if $Y=\beta_0+\beta_1Z+\eta$ and
$\varepsilon=\widehat Y-Y$, the plug-in slope error contains
$\operatorname{Cov}(Z,\varepsilon)/\operatorname{Var}(Z)$. Small average
prediction error can therefore coexist with materially biased subgroup or
treatment-effect estimates. This paper evaluates a predictor; it neither tests
moment preservation for an arbitrary downstream analysis nor supplies the
human-validation-sample correction required for calibrated inference.

\section{A latent-trait account of digital twins}
Treat respondent $i$ as a conditional response distribution
$P_i(y\given q)$, where $q$ is a survey stimulus and $y$ is a labeled option.
A digital twin is the approximation induced by model $M$ from observed
respondent context $c_i$:
\[
  \twin_i(y\given q,c_i)=M(y\given q,c_i).
\]
The object being validated is a conditional distribution, not a claim that a
prompt reproduces the respondent's complete psychology. Over a target-question
distribution $\mathcal Q$, ideal fidelity would minimize
\[
 D_i=\E_{q\sim\mathcal Q}\!
 \left[\KL\{P_i(\cdot\given q)\,\|\,\twin_i(\cdot\given q,c_i)\}\right].
\]
Only one realized answer is observed for each respondent--question pair, so the
study estimates proper-scoring-rule risk averaged over respondents and held-out
questions rather than estimating $D_i$ separately for each person.

A useful conceptual decomposition distinguishes an \emph{encoder} from a
\emph{decoder}. Let $z_i$ denote latent traits relevant to survey behavior.
The encoder $\widehat P_M(z\given c_i)$ infers those traits from demographics
and observed answers. The shared decoder $\widehat P_M(y\given q,z)$ maps a
trait configuration and new question to an answer distribution:
\[
 \twin_i(y\given q,c_i)=
 \int \widehat P_M(y\given q,z)\widehat P_M(z\given c_i)\,dz.
\]
Same-battery answers can help the encoder locate a respondent in a low-
dimensional attitude or behavioral space, while pretrained semantic knowledge
helps the decoder relate that space to new wording and option labels. Failure
can arise in either component: thin context leaves the encoder uncertain, and
an unfamiliar or weakly related target can defeat decoder transfer.

The same decomposition separates aggregate level from individual
heterogeneity. If $\widehat p_{ij}$ is respondent $i$'s predicted option vector
and $\overline p_j=n_j^{-1}\sum_i\widehat p_{ij}$, then
\[
 \widehat p_{ij}=\overline p_j+
 (\widehat p_{ij}-\overline p_j).
\]
A method may estimate $\overline p_j$ well while assigning deviations to the
wrong people. Aggregate marginal fit therefore validates the first component;
held-out respondent-level NLL and Brier validate the complete prediction.

\section{Related literature}
Language-model simulations have been proposed as synthetic samples and
economic agents \citep{argyle2023out,horton2023homo}. That work motivates the
possibility that pretrained models encode population-level regularities, but a
simulated sample is not automatically a twin of observed individuals. Recent
large-scale twin datasets instead condition on extensive respondent records
and evaluate held-out behaviors, emphasizing both predictive gains and the
remaining gap from perfect respondent copies \citep{toubia2025twin}.

This paper studies the limited-record setting: a small number of labeled
answers and readable covariates must support zero-shot transfer to a new
question. Work on survey simulation with limited human data likewise stresses
that prompting, fine-tuning, and statistical rectification answer different
parts of this problem \citep{krsteski2025valid}. Our contribution is an exact-
information evaluation framework and the \zwill{} implementation that makes such
comparisons reproducible across public microdata batteries. The LLM and
statistical benchmark receive identical case-specific evidence, while the
serialized Expected Parrot job/result boundary keeps model execution auditable
and portable across providers.

The paper is also related to direct one-shot forecasts of survey marginals and
their use as priors for survey estimation
\citep{horton2026marginalpredictors,horton2026surveypriors}. A marginal forecast
asks what the population will say without representing particular respondents.
A digital twin makes the stronger claim that respondent context predicts who
will select each option. Marginal recovery is therefore an important downstream
property here, but respondent-level proper scores remain the validation target.

Most directly, \citet{hullman2026human} argue that resemblance on observed
tasks is partial evidence of fidelity, not a statistical warrant for replacing
humans on unobserved tasks. They contrast such heuristic validation with
calibration methods that retain human labels and explicitly correct surrogate
error. \citet{ludwig2025llm} formalize the distinction. For prediction, risk
generalization requires a defensible scenario population and protection
against training leakage. For downstream estimation, LLM measurements must
preserve the moments identifying the target parameter or be combined with a
human validation sample. Our contribution sits deliberately on the prediction
side: it makes the finite benchmark and its information contract credible. It
does not convert predictive fidelity into a general interchangeability claim.

This distinction also clarifies what \zwill{} presently does and does not do.
Its leakage audit detects target information in a constructed persona; it
cannot inspect a proprietary model's pretraining corpus. Its bootstrap compares
methods over observed respondents; it does not manufacture a probability
sample of survey scenarios. Its exports could supply jointly labeled human and
model records to a rectification or prediction-powered-inference procedure,
but the current report does not implement such a downstream estimator. These
are separate contracts that software should label rather than blur.

\section{Data and cohort}
We include every normalized public-microdata battery in the study corpus with
at least two complete-case respondents.
Two batteries are excluded before analysis: one has no complete cases and one
has only a single complete case. The retained surveys span spirituality,
local-news sources, social-media use, attitudes toward AI, gender comparisons,
labor-market skills, climate policy, views of social media, Asian-American
origin-group attitudes, remote work, wellbeing, mobility and fairness,
Brazilian digital credit, municipal vignettes, and democratic norms.

The normalized source files contain human-readable question wording, option
labels, respondent weights, and respondent-level answer matrices. The
converter treats code expansion as a validation boundary: source codes are
mapped to labels before import, unmapped values become explicit missing values,
and raw metadata and respondent files are registered as provenance artifacts.
Pew profile codes for age, gender, education, race/ethnicity, party/ideology,
income tier, region, and metropolitan status are expanded where available.
Opaque covariate codes are omitted from prompts rather than presented as
personality evidence.

\begin{table}[ht]
\centering
\scriptsize
\resizebox{\textwidth}{!}{\begin{tabular}{lrrrr}\toprule
Battery & Items & Respondents & Held out & Predictions \\
\midrule
W132\_SPIRESSENT1 & 10 & 8,593 & 5 & 1,000 \\
W141\_PROVIDER & 9 & 5,036 & 5 & 1,000 \\
W150\_SMUSE & 12 & 9,749 & 5 & 1,000 \\
W152\_HUMANVAI & 8 & 5,318 & 5 & 1,000 \\
W154\_DIFF1 & 5 & 6,104 & 3 & 600 \\
W157\_SKILLIMP & 9 & 5,333 & 5 & 1,000 \\
W158\_CCPOLICY & 6 & 9,214 & 4 & 800 \\
W163\_SM9 & 5 & 5,020 & 3 & 600 \\
ASIAN2022\_23\_ORIGINBATTERY & 9 & 3,357 & 5 & 1,000 \\
GALLUP\_REMOTE\_COVID\_gallup\_remote\_covid & 9 & 1,878 & 5 & 1,000 \\
GALLUP\_WELLBEING\_gallup\_wellbeing & 10 & 5,915 & 5 & 1,000 \\
DATAVERSE\_478TLG\_KOREA\_MOBILITY\_FAIRNESS & 18 & 24,777 & 5 & 1,000 \\
DATAVERSE\_6RK7RD\_BRAZIL\_DIGITAL\_CREDIT & 19 & 647 & 5 & 1,000 \\
DATAVERSE\_DNZT11\_DNZT11\_VIGNETTE\_OUTCOMES & 6 & 1,199 & 4 & 800 \\
DATAVERSE\_IQ6EBZ\_IQ6EBZ\_2024\_DEMOCRACY\_NORMS & 15 & 1,798 & 5 & 1,000 \\
\bottomrule\end{tabular}
}
\caption{Executed study cohort. Predictions equal sampled respondents times
held-out questions for one construction arm and one model.}
\label{tab:cohort}
\end{table}

The candidate funnel is determined before inference. Table~\ref{tab:exclusions}
lists every normalized battery that failed the minimum complete-case rule; no
battery was removed because of predictive performance.

\begin{table}[ht]
\centering
\small
\begin{tabular}{p{0.42\textwidth}p{0.50\textwidth}}\toprule
Candidate battery & Pre-inference exclusion reason \\
\midrule
Trust and corruption (Harvard Dataverse) & No complete-case respondents in the normalized source. \\
Los Angeles police-stop experience (Harvard Dataverse) & Only one complete-case respondent; respondent-level validation is not meaningful. \\
\bottomrule\end{tabular}

\caption{Candidate normalized batteries excluded before plan approval and
remote inference. Reasons are read from the \zwill{} study manifest generated by
the ingestion pipeline.}
\label{tab:exclusions}
\end{table}

\section{Executed design}
\subsection{Sampling and held-out targets}
Within each battery, we hold out up to five items spanning the source order.
Because a multi-target plan removes every target from the shared context pool,
we retain at least two non-target battery items. This yields three held-out
targets for five-item batteries, four for six-item batteries, and five for
larger batteries. We sample 200 complete-case respondents per survey using
seed 20260730. Sampling and targets were stored in approved
\texttt{twin-experiment} plans before job export. The resulting primary design
contained 13,800 predictions:
\[
 \sum_s n_s \times J_s \times 1\text{ approach}\times1\text{ model}.
\]
We then executed a second, by-question design over the same 69 targets. Each
target received its own approved plan and 200-prediction job. Only that target
was hidden; up to eight other answers, including answers that were targets in
the simultaneous design, could enter the prompt. This added another 13,800
forecasts. The two designs therefore distinguish a fixed twin constructed
without the whole target set from a more informed twin reconstructed for each
question.

\subsection{Digital-twin construction}
For each respondent--target pair, the model receives:
\begin{enumerate}[nosep]
  \item the survey and held-out question wording;
  \item the complete list of labeled held-out response options;
  \item at most eight observed answers to non-held-out battery questions; and
  \item readable respondent profile metadata when the source supplied it.
\end{enumerate}
The prompt asks the model to act as the described respondent and return only a
JSON object containing one probability for every option plus a brief note.
Probabilities are kept as soft forecasts; no sampled model ``answer'' is used
for primary scoring. The construction uses \texttt{openai:gpt-5.5} with the
provider defaults exported by EDSL; no temperature parameter was supplied.

The prompt never contains the held-out answer, its observed sample or
population marginal, respondent weights, target-specific downstream
variables, or raw respondent identifiers as substantive evidence. The EDSL job
package does retain the actual answer as hidden scenario metadata so
\zwill{} can score it after inference; the rendered user prompt does not
expose that field.

Construction is performed anew for every respondent--target case rather than
by creating an unconstrained conversational agent. \zwill{} first resolves the
approved respondent sample and target-specific exclusion set. It then selects
up to eight eligible observed answers in stable survey order, attaches only
metadata fields that have readable labels, and renders a self-contained EDSL
scenario. Conceptually, the user prompt has four blocks:
\begin{quote}\small
\texttt{Respondent profile: [readable covariates]}\\
\texttt{Observed answers: [question text $\rightarrow$ labeled answer]}\\
\texttt{Question to predict: [held-out wording and all options]}\\
\texttt{Return: \{"probabilities": \{option: probability, ...\}, "note": ...\}}
\end{quote}
The output contract requires every declared option exactly once, numeric
probabilities between zero and one, and a total of one. Import records both the
rendered \texttt{observed\_answers} array and hidden actual answer. This stored
array is important: validation does not attempt to reconstruct what the model
might have seen from global survey settings; it gives XGBoost precisely that
case-specific array.

The LLM is therefore doing structured analogical prediction. It can connect
the semantics of the new item and its options to a respondent's other stated
positions, while the probability contract preserves uncertainty. No fine
tuning, retrieval from other respondents, empirical target marginal, or
question-specific demonstration is supplied. A new target can be forecast as
soon as its wording and option labels are available because construction does
not estimate parameters for that target.

\subsection{What a prediction task looked like}
The targets were ordinary closed-ended survey questions rather than abstract
benchmark labels. Illustrative examples (lightly shortened here) are:
\begin{itemize}[leftmargin=1.5em,itemsep=2pt,topsep=3pt]
  \item \textbf{Social-media use:} ``Do you ever use Facebook?'' (Yes; No).
  \item \textbf{Gender comparisons:} ``Are men and women basically similar or
  different in their hobbies and personal interests?'' (Similar; Different).
  \item \textbf{Origin-group attitudes:} ``What is your opinion of the U.S.?''
  (Very favorable through Very unfavorable; five options).
  \item \textbf{Municipal vignettes:} ``Based on Vignette 1, are city leaders
  to blame for trash and recycling bins not being picked up?'' (Strongly
  disagree through Strongly agree; seven options).
  \item \textbf{Democratic norms:} ``How confident are you that all votes will
  be counted accurately?'' (No confidence through Very much confidence; four
  options).
\end{itemize}
The model jobs used the complete imported wording and labeled option lists.

For a concrete climate-policy case, the held-out question asked whether the
respondent favored planting roughly one trillion trees to absorb carbon
emissions. The prompt described a 30--49 year-old woman with postgraduate
education, moderate ideology, and a moderate/liberal Republican affiliation.
It also reported that she favored requiring oil and gas companies to seal
methane leaks and favored home-energy-efficiency tax credits. The prompt then
listed the target options \emph{Favor} and \emph{Oppose} and required a JSON
probability vector. GPT-5.5 returned 0.90 and 0.10, respectively; the held-out
answer was Favor. This is one successful row, not a cherry-picked performance
estimate: the reported metrics aggregate all 200 respondents and all
predeclared targets in the battery.

The example also makes the information structure tangible. The twin is not
asked to answer from demographics alone. It combines a readable profile with
the respondent's answers to neighboring, non-held-out items. In a six-item
battery with four shared targets, only the two remaining items are eligible
context; in larger batteries as many as eight are shown. At scoring time the
same respondent's hidden target answer is joined back to the forecast.

\subsection{Approval, execution, and preservation}
Each survey was imported and committed before model work. We rendered a
pre-run survey report, registered the construction approach, initialized and
validated a plan, explicitly approved the plan, and compared approved and
exported counts. The 15 simultaneous exports and 69 by-question exports each
matched their approved counts exactly: 13,800 predictions per design. Expected
Parrot executed the serialized EDSL packages
and returned immutable Results packages. \zwill{} preserved each raw
package and its hash before extracting probability rows. A later leakage
correction added 200 predictions, bringing the executed total to
\TotalPredictions.

For model robustness, we predeclared five substantively varied batteries
(gender comparisons, climate policy, origin-group attitudes, Brazilian
digital credit, and democratic norms) containing 22 targets. We reran the
by-question construction for all 200 respondents with Claude Sonnet 4.6 and
Gemini 3.5 Flash, adding $22\times200\times2=8{,}800$ forecasts. Model-specific
approach identifiers keep provider provenance in the job identity while
leaving the prompt information policy unchanged. The matching GPT-5.5 rows
were already present. All three models are evaluated on identical cases and
against one shared exact-context XGBoost baseline.

\section{A reproducible \zwill{} workflow}
\label{sec:tutorial}
This section presents the study as a worked \zwill{} tutorial. The code
fragments are shortened versions of the executable Makefile recipes. A
declarative study manifest supplies source-specific metadata, while \zwill{}
performs batch conversion, manages survey and study state, and conducts
validation. Expected Parrot's \texttt{ep} command executes exported model jobs.
This separation is intentional: creating a survey or plan does not make
a remote model call. All commands below run from an isolated paper work
directory. After every state-changing stage, \texttt{zwill next --survey NAME}
reports the current state and expected next command.

\subsection{Initialize, preserve provenance, and import labeled data}
A project begins with an explicit project and a declarative matrix specification.
The specification names identifier and weight columns, maps raw option and
profile codes to readable labels, declares missing codes, and identifies each
item's source and data variables. Conversion completes before project state is
changed.

\begin{lstlisting}[style=shell]
zwill init
zwill project create microdata_twins_v2 --use
\end{lstlisting}
These commands create an isolated project; they do not run a model.

\begin{lstlisting}[style=shell]
zwill study init --path ingestion_manifest.json
\end{lstlisting}
This creates only a minimal, empty declaration---a schema version, an empty
defaults object, and an empty survey list. It neither searches the filesystem
nor guesses which files belong together.

\begin{lstlisting}[style=shell]
zwill study set-matrix-defaults \
  --manifest ingestion_manifest.json \
  --covariate-codebook config/covariate_codebook.json \
  --data-variable-template 'item_{item}' \
  --question-name-template '{battery}_{item}' \
  --missing-code NA --missing-code Refused
\end{lstlisting}
This records visible shared policy in the manifest. Every default came from an
argument; the command performs no inference about code meanings or naming.

\begin{lstlisting}[style=shell]
zwill study add-matrix --manifest ingestion_manifest.json \
  --survey w154_diff1 --source-id W154_DIFF1 \
  --metadata data/W154_DIFF1_metadata.json \
  --respondents data/W154_DIFF1_respondents.csv
\end{lstlisting}
One invocation declares one source pair. Repeating it builds the 15-battery
list and rejects duplicate survey identifiers. The resulting JSON is ordinary,
reviewable configuration; for this example its essential structure is:
\begin{lstlisting}[style=json]
{
  "matrix_defaults": {
    "covariate_codebook": "config/covariate_codebook.json",
    "metadata_overlay": {
      "data_variable_template": "item_{item}",
      "question_name_template": "{battery}_{item}",
      "missing_codes": ["NA", "Refused"]
    }
  },
  "surveys": [{
    "survey": "w154_diff1", "source_id": "W154_DIFF1",
    "matrix": {
      "metadata": "data/W154_DIFF1_metadata.json",
      "respondents": "data/W154_DIFF1_respondents.csv"
    }
  }]
}
\end{lstlisting}

\begin{lstlisting}[style=shell]
zwill study check --manifest ingestion_manifest.json
\end{lstlisting}
This read-only preflight resolves every declared path and prints, by battery,
the source identifier, respondent count, question count, matrix columns, and
required columns. It checks metadata, overlays, codebooks, duplicate ids, and
column coverage without creating or changing a \zwill{} project.

\begin{lstlisting}[style=shell]
zwill study ingest --manifest ingestion_manifest.json \
  --out "$PWD/imports" --apply
\end{lstlisting}
The study manifest supplies shared matrix defaults and one record per battery;
individual records can override paths or metadata. It supports source and
survey path templates, metadata overlays, readable item-name templates, and a
shared covariate codebook. \zwill{} first converts \emph{every} matrix; only if
all conversions succeed does it apply any survey. The adapter checks required
columns and codebooks, normalizes numeric codes, labels answers and covariates,
writes a conversion manifest, archives raw sources, and imports the three
canonical JSONL streams. Unknown answer or profile codes, or invalid weights,
produce \texttt{issues.jsonl} and stop the batch before partial application.
The generated streams remain inspectable intermediate artifacts.

\begin{lstlisting}[style=shell]
zwill commit --survey w154_diff1
\end{lstlisting}
Commit freezes the labeled scoring truth after the clean conversion is reviewed.

\begin{lstlisting}[style=shell]
zwill survey report --survey w154_diff1 --format html \
  --path reports/w154_diff1/survey-profile.html
\end{lstlisting}
The report exposes wording, options, distributions, and missingness before paid
inference. The executable workflow runs \texttt{zwill next} after each stage.

\subsection{Register the construction and predeclare the experiment}
The construction specifies what information may create a twin. The following
sequence registers complete-case sampling, at most eight context answers,
readable profile metadata (included by default), and GPT-5.5. A shared study
manifest then writes all 15 primary plans and their projected call counts.

\begin{lstlisting}[style=shell]
zwill twin-approach add --survey w154_diff1 \
  --approach-id battery_context \
  --name "Held-out battery context" \
  --description "Predict targets from non-target answers and profile fields." \
  --complete-cases --context-question-count 8 \
  --model openai:gpt-5.5
\end{lstlisting}
The approach is a reusable information policy rather than a particular target.

\begin{lstlisting}[style=shell]
zwill study prepare --manifest study_manifest.json \
  --mode simultaneous --out "$PWD/workflow" \
  --results-dir "$PWD/results" --approved-by "Reviewer Name"
\end{lstlisting}
The manifest holds shared defaults plus each survey's source identifier,
targets, stable plan identifier, and target-specific leakage exclusions.
Preparation composes expansion, recorded approval, and export. It removes
hand-written shell fan-out while retaining plans in \texttt{workflow/plans} and
packages in \texttt{workflow/jobs}. Existing plans are never overwritten,
preserving approvals and review notes on resumed builds.

\begin{lstlisting}[style=shell]
zwill twin-experiment validate \
  --input-path plans/w154_diff1_validation_v1.json
\end{lstlisting}
The plan materializes targets, respondent sample, seed, metric, leakage policy,
and expected prediction count. Validation does not approve spending or execute
inference.

The plural \texttt{--heldout-questions} argument matters. All three W154
targets are removed from context together. Each respondent therefore has the
same profile and observed non-target answers when predicting A, B, and C; only
the requested target and its options change. Here D and E remain as context.
This is a fixed-information, simultaneous multi-target holdout, not a
round-robin in which B and C become visible while A is predicted. Its planned
call count is $200\times3\times1\times1=600$.

The by-question benchmark instead passes \texttt{--mode by-question} to
\texttt{zwill study prepare}, creating three separate W154
plans. When A is the sole target,
B--E are eligible context; when B is the target, A and C--E are eligible. This
requires new twin inference because the rendered personas differ from those in
the simultaneous job.

The plan JSON records targets, approach, sample size, seed, model, primary
metric, approval state, and expected call count. Separating this object from
the provider package allows scientific review before an executable job exists.

\subsection{Review, export, and execute through Expected Parrot}
Approval is the spending boundary. Export turns the approved design into an
EDSL \texttt{Jobs.ep} package but still performs no inference. The external
\texttt{ep run} command submits that package and preserves a raw
\texttt{Results.ep} package.

The \texttt{prepare} command records the reviewer on every expanded plan, then
validates the approved designs, renders scenarios, verifies counts, and writes
portable \texttt{Jobs.ep} packages plus one resumable
\texttt{run-manifest.json}. Each run record contains its job path, expected
result path, survey, plan, model, and exact \texttt{ep} command.

\begin{lstlisting}[style=shell]
ep run-manifest workflow/jobs/run-manifest.json \
  --parallel 3 --execute
\end{lstlisting}
Only this command can perform remote calls, and \texttt{--execute} is the explicit
spending gate. Without it the command only verifies cached results. It compares
each result count with the manifest, skips complete packages, and preserves an
incomplete package under a \texttt{partial} name before resuming it.

The export manifest compares approved and exported counts; a discrepancy
should stop execution and trigger reapproval. Thus an interrupted build resumes
without purchasing completed predictions again.

\subsection{Import results and audit the actual prompts}
The importer joins responses to hidden scenario metadata, checks the
probability contract, and stores parsed predictions. Before looking at scores,
the run report exposes the template, rendered examples, response text,
construction metadata, and parse issues.

\begin{lstlisting}[style=shell]
zwill study finish --manifest workflow/jobs/run-manifest.json \
  --results-dir results --replace \
  --audit-out "$PWD/reports/run-audits" \
  --out "$PWD/reports/{survey}/validation-exact" \
  --embedder edsl --embedding-model text-embedding-3-small \
  --seed 20260730
\end{lstlisting}
Finishing composes import, prompt audit, and validation. It refuses to audit or
score incomplete results. The importer verifies scenario metadata and the
probability contract while retaining raw responses; the reports then expose
rendered personas, hidden targets, parameters, responses, and parse issues.

The audit must show the permitted profile and non-target answers, but not the
target answer, target marginal, respondent weight, or target-revealing
downstream fields. Successful parsing alone is not validity evidence: a
perfectly structured prompt can still leak its answer.

\subsection{Fit the fair baseline and build the evidence bundle}
The validation stage inside \texttt{study finish} performs the decisive comparison. It audits
context--target associations, fits the conditional baseline, computes paired
bootstrap intervals, and renders machine-readable and visual diagnostics.

This groups imported job ids by battery and invokes the exact same validation
contract for each. It audits leakage, fits matched XGBoost, bootstraps paired
score differences, and writes one evidence bundle per battery. A missing
conditional baseline fails validation.

\begin{lstlisting}[style=shell]
zwill report build --survey w154_diff1 \
  --jobs f2e86020755065ce \
  --path reports/w154_diff1/final --seed 20260730
\end{lstlisting}
This gathers survey, prediction, audit, and validation artifacts into a linked
final-report bundle.

\begin{lstlisting}[style=shell]
zwill report render --survey w154_diff1 \
  --jobs f2e86020755065ce \
  --path reports/w154_diff1/final --seed 20260730
\end{lstlisting}
Rendering turns that bundle into its human-readable HTML readout.

The study validator requires its conditional baseline rather than allowing a
weaker twin-versus-marginal report to appear complete. In this paper, baseline
training and scoring are restricted to the same sampled
respondent identifiers and the exact per-case \texttt{observed\_answers} list
preserved by the twin importer. Thus the cap, missingness pattern, and any
target-specific exclusion also match. Forcing the \texttt{edsl} embedder
prevents silent fallback to a different representation.

The validation directory contains the report, baseline predictions, leakage
audit, bootstrap estimates, confidence-interval graphic, calibration graphic,
and a manifest. Those artifacts, not hand-copied numbers, feed the paper's
cross-survey summary, tables, and plot.

\subsection{Add battery-local one-shots and named model comparisons}
The marginal elicitation is deliberately partitioned by battery. Adding a new
battery creates one new job rather than invalidating a monolithic job or
repurchasing completed forecasts.

\begin{lstlisting}[style=shell]
zwill study one-shot-build --manifest study_manifest.json \
  --out "$PWD/jobs/one-shot" --model openai:gpt-5.5
\end{lstlisting}
This writes one no-persona probability \texttt{Jobs.ep} package per battery and
a run manifest containing the corresponding remote commands and result paths.
Each question is asked once for its population distribution.

\begin{lstlisting}[style=shell]
zwill study one-shot-import \
  --manifest jobs/one-shot/run-manifest.json \
  --results-dir results/one-shot \
  --path reports/one-shot/probability-jobs.json --replace
\end{lstlisting}
After the recorded jobs are executed by \texttt{ep run}, import produces a
stable battery-to-job mapping. Passing that mapping to \texttt{study validate}
scores the one-shot distribution on exactly the respondent--question cases used
by the twins, without further model calls.

The same study manifest declares the provider robustness exercise as a named
comparison: its five-battery subset, models, approach identifiers, mode, and
plan-id template are data rather than shell branches.

\begin{lstlisting}[style=shell]
zwill study expand --manifest study_manifest.json \
  --comparison frontier_models \
  --out "$PWD/plans/model-comparison"
\end{lstlisting}
The result is the full survey--question--model plan matrix. The same study-level
approve, export, import, audit, and validate commands then operate on its
manifest. Validation can include the primary run manifest so every provider is
compared with the identical stored XGBoost arm and cases.

\subsection{Normalize cross-survey evidence and realized costs}
Paper-specific plotting code should not reopen version-dependent internal
bundle files. \zwill{} first normalizes every validation bundle into stable,
tidy study artifacts.

\begin{lstlisting}[style=shell]
zwill study summarize --manifest study_manifest.json \
  --bundle-template 'zwill_work/reports/{survey}/validation-exact' \
  --out "$PWD/study/simultaneous"
\end{lstlisting}
This writes \texttt{surveys.csv}, \texttt{questions.csv},
\texttt{marginals.csv}, \texttt{options.csv}, \texttt{summary.json}, and
\texttt{aggregates.json}. They
contain explicit uniform, empirical-marginal, conditional-XGBoost, and one-shot
comparisons, including paired intervals where available. NLL deltas consistently
mean model loss minus benchmark loss, so negative values favor the twin.
\texttt{aggregates.json} owns deterministic cross-question bootstrap summaries,
NLL and Brier benchmark wins, clean-claim gate counts, calibration summaries,
and marginal MAE and correlation. The paper does not recompute those statistics.

\begin{lstlisting}[style=shell]
zwill study plot --summary-dir "$PWD/study/by-question" \
  --out "$PWD/study/plots" --model openai:gpt-5.5
\end{lstlisting}
This creates reusable dependency-free SVG diagnostics directly from the stable
tables. The paper's remaining Python is presentation-specific: it combines
designs, applies journal styling, and emits LaTeX, but does not reopen internal
validation bundles or reconstruct \zwill{} scores.

\begin{lstlisting}[style=shell]
zwill study costs --manifest study_manifest.json \
  --path "$PWD/study/costs.json"
\end{lstlisting}
Result-package globs are declared relative to the manifest. The command
deduplicates matches and records realized call counts, costs, per-model totals,
and a per-file audit trail; an empty match is an error rather than a silent zero.

\subsection{Reproduce the complete paper}
The commands above expose each reproducible boundary. The Makefile connects
their manifests and stamp files; it does not encode battery-specific analysis
or model matrices. It applies the chain to all 15 batteries and then generates
the paper:

\begin{lstlisting}[style=shell]
cd papers/zwill_microdata_twins
make reproduce ALLOW_REMOTE=1 APPROVED_BY="Reviewer Name"
\end{lstlisting}

Without \texttt{ALLOW\_REMOTE=1}, a missing result stops the build at
inference. With cached result and validation packages, \texttt{make all} runs
only deterministic summarization, plotting, and typesetting. This makes the
common operation safe while retaining a genuine source-to-PDF replication
target for an authorized operator.

\section{Evaluation}
\subsection{Outcomes}
The primary metric is respondent-level negative log likelihood, with Brier
score, probability assigned to the observed answer, calibration, and top-1
accuracy as secondary diagnostics. For respondent $i$ with observed option
$y_{ij}$, $K_j$ options, forecast vector $\widehat p_{ij}$, and one-hot outcome
$e_{ij}$, the two proper scores are
\[
 \operatorname{NLL}_{ij}=-\log\widehat p_{ij,y_{ij}},
 \qquad
 \operatorname{Brier}_{ij}=\sum_{a=1}^{K_j}
 (\widehat p_{ija}-e_{ija})^2.
\]
Lower values are better. Top-1 accuracy and mean probability assigned to the
actual answer are useful diagnostics, but neither is a proper scoring rule;
accuracy in particular discards forecast uncertainty.

Calibration asks whether events assigned probability $p$ occur approximately
$p$ of the time. \zwill{} reports expected calibration error over probability
bins and exposes overconfident misses for prompt-level inspection. We treat
calibration as a diagnostic rather than allowing it to override the paired NLL
gate. The uniform reference assigns $1/K_j$ to every option. The empirical
marginal reference assigns each respondent the target distribution observed in
the scored sample; because that distribution requires target answers, it is an
oracle reference rather than a deployable new-question method.

At the aggregate level, averaging respondent forecasts gives
$\widehat\pi_{ja}=n_j^{-1}\sum_i\widehat p_{ija}$. We compare it with the
observed option share $\pi_{ja}$ using option-share MAE, Jensen--Shannon
divergence, modal-option agreement, and rank correlation. These quantities test
level and ordering. Respondent-level NLL and Brier test whether probability is
assigned to the correct individuals. Finally, \zwill's joint-structure
diagnostics compare predicted and observed cross-question association; this
guards against twins that reproduce each marginal while washing out dependence
between answers.

\subsection{Conditional baseline}
The comparison baseline is an exact-information-set XGBoost classifier. For
each target, \zwill{} reads the \texttt{observed\_answers} stored with
every imported twin prediction. Those question names define that respondent's
permitted baseline context. XGBoost receives neither additional respondents
nor additional answers. Consequently, the simultaneous design excludes the
whole target set and respects the eight-answer cap, whereas the by-question
design may use the other targets only because the corresponding twin prompt
also used them. Target-specific exclusions and respondent-level missingness
carry through automatically.

More precisely, the baseline converts each training question into one row per
respondent--candidate-option pair. The binary outcome is one when the candidate
option equals that respondent's observed answer and zero otherwise. A row
contains (i) the question-text vector, (ii) the candidate-option vector,
(iii) one-hot respondent covariates, (iv) cosine-similarity features comparing
the candidate with the average text representation of the respondent's other
answers, and (v) the candidate's position in the option list. When predicting
target $j$, training labels come only from questions present in the twin's
permitted context. A training row for context question $k$ excludes $k$ itself
from its respondent profile, as well as every answer not permitted for the
target case. Thus $j$ cannot enter indirectly as a feature on another
question's training row.

Formally, let $C_{ij}$ be the exact set of context answers shown to respondent
$i$'s twin when target $j$ is predicted, and let $e(\cdot)$ denote the fixed
text embedding.  For a training context question $k\in C_{ij}$ and candidate
option $a\in\mathcal A_k$, define
\[
 z_{ika}^{(j)}=
 \left[
 e(q_k),\ e(a),\ x_i,\
 \cos\!\left(e(a),\frac{1}{|C_{ij}\setminus\{k\}|}
       \sum_{\ell\in C_{ij}\setminus\{k\}}e(y_{i\ell})\right),\
 \operatorname{pos}(a,k)
 \right],
 \qquad
 r_{ika}=\mathbf 1\{y_{ik}=a\},
 \]
where $x_i$ is the one-hot respondent profile. The exclusion of $k$ from the
profile prevents the training label from reappearing as an input; because
$C_{ij}$ comes directly from the saved twin prompt, no unshown response can
enter the baseline.

XGBoost estimates an additive tree score
\[
 F_M(z)=\sum_{m=1}^{M}\eta f_m(z), \qquad f_m\in\mathcal T_6,
\]
by minimizing regularized binary log loss over the respondent--question--option
rows,
\[
 \widehat F_j=\arg\min_F
 \sum_i\sum_{k\in C_{ij}}\sum_{a\in\mathcal A_k}
 \left[-r_{ika}\log\sigma(F(z_{ika}^{(j)}))
 -(1-r_{ika})\log\{1-\sigma(F(z_{ika}^{(j)}))\}\right]
 +\Omega(F).
\]
We use $M=300$ trees, maximum depth 6, learning rate $\eta=0.1$, row and
column subsampling of 0.8, L2 regularization of 1, histogram tree construction,
and log-loss evaluation. For target $j$, each option receives a score and the
reported categorical forecast is
\[
 \widehat p_{ija}=
 \frac{\exp\{\widehat F_j(z_{ija}^{(j)})\}}
      {\sum_{b\in\mathcal A_j}\exp\{\widehat F_j(z_{ijb}^{(j)})\}}.
\]

This setup can score a genuinely new question because it does not require a
question-specific coefficient or any labeled examples for that target. At
prediction time, the new question and each of its options are embedded using
the same transformation. The model scores each candidate from text features,
respondent covariates, and similarity to the respondent's other-answer
profile; candidate scores are normalized to sum to one. Transfer is possible
because the tree learns cross-question regularities---for example, whether a
candidate semantically resembles options previously selected by a respondent,
or whether a covariate predicts response tendencies across items. The baseline
job identifier includes a fingerprint of all respondent--target context lists,
so the simultaneous and by-question predictions cannot overwrite one another.

All fair baselines use OpenAI's \texttt{text-embedding-3-small}, routed through
Expected Parrot, for question text, option text, and respondent other-answer
profiles. A single semantic representation is held fixed across all 15
batteries. The embedding model supplies features only;
it is not trained on these surveys and does not receive the held-out answers.

\subsection{Exploratory semantic-proximity analysis}
We ask whether a target is easier when its wording is semantically closer to
the questions actually shown in the twin's context. Using the same
\texttt{text-embedding-3-small} representation, let $u_j$ denote the normalized
embedding of target question $j$, and let $C_{ij}$ be the context question names
stored with respondent $i$'s imported prediction. We define
\[
 s_{ij}=\frac{1}{|C_{ij}|}\sum_{k\in C_{ij}}u_j'u_k,
 \qquad s_j=\frac{1}{n_j}\sum_i s_{ij}.
\]
This uses the context that was actually rendered, so missing answers, context
caps, and target-specific exclusions are respected. We also compute the mean,
over respondents, of the maximum target--context similarity as an alternative
``nearest question'' measure.

At the question level we estimate
\[
 Y_{jb}=\alpha_b+\beta s_{jb}+\varepsilon_{jb},
\]
where $alpha_b$ is a battery fixed effect. Outcomes are twin mean NLL,
accuracy, and twin advantage over XGBoost, defined as
$-(\mathrm{NLL}_{twin}-\mathrm{NLL}_{XGB})$ so that larger values mean a larger
twin advantage. Confidence intervals resample the 20 batteries with replacement
and refit the fixed-effect slope 10,000 times. This is exploratory: 69
questions were available in the primary collection and 107 after the GSS
extension; proximity is not randomly assigned, and the analysis
cannot distinguish semantic similarity from unmeasured question difficulty.
For visualization, we additionally fit a local-linear LOESS curve with a 0.75
span to the within-battery residuals. Its pointwise 95\% band comes from 2,000
of the same whole-battery bootstrap resamples. The linear coefficient remains
the reported summary; the smoother is a diagnostic for nonlinearity and
leverage rather than a second confirmatory specification.

\begin{lstlisting}[style=shell]
zwill study semantic-proximity \
  --manifest jobs/by_question/run-manifest.json \
  --manifest jobs/gss2024_eligible200/run-manifest.json \
  --questions study/by-question/questions.csv \
  --questions cross_survey/gss2024/questions.csv \
  --out study/semantic-proximity --model openai:gpt-5.5 \
  --embedder edsl --embedding-model text-embedding-3-small
\end{lstlisting}
The command writes the question-level design matrix, within-battery residuals,
fixed-effect estimates, and battery-bootstrap intervals. The Makefile treats
those artifacts as inputs to the publication plot.

\subsection{Uncertainty and decision rule}
Each bundle uses 1,000 paired bootstrap resamples over respondents and a 95\%
confidence interval. We define
\[
 \Delta_{\mathrm{NLL}} =
 \mathrm{NLL}_{\mathrm{twin}}-\mathrm{NLL}_{\mathrm{conditional}}.
\]
A survey passes only if the entire interval is below zero, fails if the entire
interval is above zero, and is inconclusive otherwise. A positive result also
requires a clean leakage audit. We do not use an accuracy improvement to
override an inconclusive or adverse NLL result.

\subsection{Prompt and leakage audits}
Before performance interpretation, we generated a run report for every job
containing the prompt template, rendered scenarios, observed context, raw model
responses, model parameters, parse counts, and import issues. The statistical
leakage audit screens context--target pairs using bias-corrected Cramér's V,
flagging values above 0.7 with at least 30 co-answered respondents. A flagged
pair blocks an unqualified positive claim even when predictive metrics look
strong.

\section{Results}
\subsection{Execution integrity and cost}
The model produced \TotalPredictions\ held-out probability predictions,
including a 200-prediction corrective run, at an observed inference cost of
\TotalCost. All \TotalPredictions\ model calls returned result rows and all probability
objects parsed successfully; there were no malformed or silently dropped rows.
The total includes 13,800 simultaneous forecasts, 13,800 newly purchased
by-question forecasts, the 200-call Gallup correction, and 8,800 forecasts in
the two-model robustness extension.

\subsection{Survey-level validation}
Figure~\ref{fig:nll-change} plots paired macro-average NLL differences relative
to the conditional baseline and their 95\% bootstrap intervals. Negative
values favor the twin. Circles show simultaneous holdout and triangles show
by-question holdout. Hollow red points mark the Gallup remote-work battery,
which is shown for transparency but excluded from clean-survey claims because
of its leakage flag. All 15 batteries have exact-context conditional baselines
under both designs. Table~\ref{tab:validation} gives the corresponding
numeric values and audit status.

\begin{figure}[ht]
\centering
\includegraphics[width=0.92\textwidth]{../output/figures/nll_change.pdf}
\caption{Change in negative log likelihood relative to XGBoost trained on the
same respondents and exact prompt-level answer context. Points are survey-level
macro averages and bars are paired 95\% bootstrap intervals. Values below zero
favor GPT-5.5. Circles denote simultaneous target holdout; triangles denote
one-target-at-a-time holdout. Hollow remote-work points have a leakage flag.}
\label{fig:nll-change}
\end{figure}

\begin{table}[ht]
\centering
\scriptsize
\resizebox{\textwidth}{!}{\begin{tabular}{lr r l r l}\toprule
Battery & Flags & $\Delta$ NLL$_{sim}$ & 95\% CI & $\Delta$ NLL$_{byq}$ & 95\% CI \\
\midrule
ASIAN2022\_23\_ORIGINBATTERY & 0 & -0.991 & [-1.162, -0.830] & -0.824 & [-0.950, -0.693] \\
GALLUP\_REMOTE\_COVID\_gallup\_remote\_covid & 1 & -0.907 & [-0.997, -0.817] & -0.822 & [-0.943, -0.705] \\
W163\_SM9 & 0 & -0.477 & [-0.581, -0.376] & -0.482 & [-0.582, -0.388] \\
GALLUP\_WELLBEING\_gallup\_wellbeing & 0 & -0.411 & [-0.481, -0.342] & -0.399 & [-0.458, -0.344] \\
W154\_DIFF1 & 0 & -0.407 & [-0.508, -0.308] & -0.135 & [-0.212, -0.067] \\
W157\_SKILLIMP & 0 & -0.353 & [-0.435, -0.268] & -0.306 & [-0.384, -0.233] \\
DATAVERSE\_6RK7RD\_BRAZIL\_DIGITAL\_CREDIT & 0 & -0.346 & [-0.391, -0.303] & -0.496 & [-0.555, -0.440] \\
DATAVERSE\_DNZT11\_DNZT11\_VIGNETTE\_OUTCOMES & 0 & -0.317 & [-0.399, -0.236] & 0.029 & [-0.062, 0.118] \\
W150\_SMUSE & 0 & -0.263 & [-0.337, -0.190] & -0.180 & [-0.244, -0.116] \\
W158\_CCPOLICY & 0 & -0.204 & [-0.287, -0.138] & -0.203 & [-0.278, -0.126] \\
DATAVERSE\_IQ6EBZ\_IQ6EBZ\_2024\_DEMOCRACY\_NORMS & 0 & -0.187 & [-0.234, -0.144] & -0.220 & [-0.271, -0.170] \\
W152\_HUMANVAI & 0 & -0.150 & [-0.194, -0.110] & -0.244 & [-0.295, -0.191] \\
W141\_PROVIDER & 0 & -0.150 & [-0.210, -0.092] & -0.190 & [-0.250, -0.134] \\
W132\_SPIRESSENT1 & 0 & -0.144 & [-0.220, -0.059] & -0.207 & [-0.271, -0.143] \\
DATAVERSE\_478TLG\_KOREA\_MOBILITY\_FAIRNESS & 0 & -0.124 & [-0.163, -0.083] & -0.176 & [-0.217, -0.137] \\
\bottomrule\end{tabular}
}
\caption{Exact-information-set NLL comparisons under simultaneous (sim) and
by-question (byq) holdout. ``Pass'' requires the complete 95\% interval to be
below zero. The flag count is common to both designs.}
\label{tab:validation}
\end{table}

Among 14 clean surveys, simultaneous-holdout twins pass the NLL gate in
\SimNLLPass\, fail in \SimNLLFail\, and are inconclusive in
\SimNLLInconclusive. Their intervals are entirely below zero in every clean
battery, with deltas ranging from roughly $-1.01$ for origin-group attitudes to
$-0.14$ for spirituality. Brier intervals also favor the twin in all 14.

Under by-question holdout, \ByQNLLPass\ clean surveys pass, \ByQNLLFail\ fail,
and \ByQNLLInconclusive\ is inconclusive. Municipal vignettes are the lone
inconclusive case ($\Delta_{\mathrm{NLL}}=-0.005$, 95\% CI
$[-0.094,0.089]$); the other 13 intervals clear zero favorably. Brier results
give the same 13--0--1 pattern. Providing both methods with answers to the other
targets therefore narrows the advantage in some batteries, but does not
generally eliminate it.

\subsection{Extension to the 2024 General Social Survey}
We next applied the same one-target-at-a-time design to five batteries extracted
from the official 2024 GSS Release 3a: confidence in institutions, socializing
frequency, abortion circumstances, government responsibility, and national
spending. The declarative extraction specification retains reviewed question
wording, labeled response options, source-variable provenance, and readable
ballot, form, and mode covariates. Each target uses 200 respondents drawn from
the target battery's complete-case pool. The twin and conditional XGBoost both
receive every other admissible answer in that battery, so this is the same
exact-information comparison as the primary by-question design.

The extension contains \GSSQuestions{} held-out questions and \GSSCalls{} model
calls, including 38 same-model no-persona one-shot forecasts, at a recorded cost
of \GSSCost. GPT-5.5 improves on conditional XGBoost by
\GSSConditionalDelta{} NLL on average across questions (95\% battery-bootstrap
interval $[\GSSConditionalLo,\GSSConditionalHi]$), winning 37 of 38 questions
and all five battery-level comparisons. Relative to the one-shot GPT-5.5
distribution, the average difference is \GSSOneShotDelta{}
($[\GSSOneShotLo,\GSSOneShotHi]$), with wins on 36 of 38 questions and all five
batteries. The option-level implied-marginal correlation is
\GSSMarginalCorrelation{} and MAE is \GSSMarginalMAE.

Figure~\ref{fig:gss-extension} and Table~\ref{tab:gss-extension} report the
battery results. The abortion battery is displayed but excluded from the
audit-clean headline count: twelve directional context--target pairs exceed the
predeclared Cram\'er's V threshold because the closely related abortion
circumstances are strongly associated. Prompt audits confirm that the held-out
answer itself was absent. We nevertheless retain the conservative gate rather
than treating substantive redundancy as automatically harmless.

\begin{figure}[ht]
\centering
\includegraphics[width=0.88\textwidth]{../output/figures/gss_extension.pdf}
\caption{Leave-one-question-out validation in five 2024 GSS batteries. Points
are battery macro-average GPT-5.5 NLL minus exact-context XGBoost NLL; bars are
paired 95\% bootstrap intervals. The hollow abortion point is audit-flagged.}
\label{fig:gss-extension}
\end{figure}

\begin{table}[ht]
\centering
\scriptsize
\resizebox{\textwidth}{!}{\begin{tabular}{lrrrrl}\toprule
Battery & Questions & Flags & $\Delta$ NLL$_{XGB}$ & 95\% CI & $\Delta$ NLL$_{1shot}$ \\
\midrule
GSS national spending & 11 & 0 & -0.418 & [-0.467, -0.370] & -0.136 \\
GSS government responsibility & 4 & 0 & -0.335 & [-0.445, -0.223] & -0.277 \\
GSS socializing frequency & 4 & 0 & -0.318 & [-0.402, -0.230] & -0.022 \\
GSS institutional confidence & 12 & 0 & -0.311 & [-0.355, -0.266] & -0.169 \\
GSS abortion circumstances (flagged) & 7 & 12 & -0.248 & [-0.311, -0.189] & -0.358 \\
\bottomrule\end{tabular}
}
\caption{GSS extension results against exact-context XGBoost and the same-model
no-persona one-shot baseline. Negative differences favor respondent twins.}
\label{tab:gss-extension}
\end{table}

\subsection{External validation in East African agricultural microdata}
To probe performance outside prominent recurring survey series, we applied the
same workflow to the KALRO baseline from the replication package for
\emph{Digital Information Provision and Behavior Change: Lessons from Six
Experiments in East Africa} (AEA Data and Code Repository/openICPSR
E186241V1). This is a pre-treatment field-study sample of
\KALRORespondents{} Kenyan respondents rather than a nationally prominent
opinion survey. The declarative extraction retains three coherent binary
batteries: awareness of four agricultural inputs, prior use of those four
inputs, and prior use of twelve farming practices. Following the authors'
cleaning logic, structurally skipped use items are coded as no use. Readable
gender, age, county, and subcounty fields are retained as profile metadata.

The extension is itself a worked \zwill{} example. First, a declarative CSV
extraction specification maps source columns, labels, missing values, profile
fields, and the structural-skip rule into three respondent matrices:

\begin{lstlisting}[style=shell]
zwill survey extract-csv \
  --data sources/openicpsr_186241/raw/KALRO_baseline_raw.csv \
  --spec ../config/openicpsr_186241_kalro_batteries.json \
  --out sources/openicpsr_186241/extracted
\end{lstlisting}
This produces reviewable metadata JSON and respondent CSV files for each
battery. Each matrix is then validated and applied to the KALRO project; the
following is the awareness battery, with the same command repeated for input
use and farming practices:

\begin{lstlisting}[style=shell]
zwill survey ingest-matrix --survey kalro_input_awareness \
  --metadata sources/openicpsr_186241/extracted/kalro_input_awareness_metadata.json \
  --respondents sources/openicpsr_186241/extracted/kalro_input_awareness_respondents.csv \
  --out sources/openicpsr_186241/converted/kalro_input_awareness --apply
\end{lstlisting}
After inspection, \texttt{zwill commit --survey kalro\_input\_awareness}
freezes the observed scoring truth (and likewise for the other two batteries).
The study manifest names all 20 targets and shared design settings. One command
then creates, records approval for, and exports the 20 leave-one-question-out
plans while leaving every plan JSON available for inspection:

\begin{lstlisting}[style=shell]
zwill study prepare \
  --manifest ../config/openicpsr_186241_kalro_study.json \
  --mode by-question --out workflow/openicpsr_186241/full \
  --plan-suffix validation_v1 \
  --results-dir results/openicpsr_186241/full \
  --approved-by "Reviewer Name"
\end{lstlisting}
Expected Parrot executes the resulting manifest independently:

\begin{lstlisting}[style=shell]
ep run-manifest \
  workflow/openicpsr_186241/full/jobs/run-manifest.json \
  --parallel 3 --execute
\end{lstlisting}
Omitting \texttt{--execute} performs a free completeness check. The no-persona
comparator is separately exportable, so it can be added or rerun without
changing the respondent-level plans:

\begin{lstlisting}[style=shell]
zwill study one-shot-build \
  --manifest ../config/openicpsr_186241_kalro_study.json \
  --out jobs/openicpsr_186241/full_one_shot \
  --model openai:gpt-5.5
\end{lstlisting}
After those three battery-level jobs run, one command imports them and writes a
survey-to-job mapping rather than coupling their lifecycle to the twin run:

\begin{lstlisting}[style=shell]
zwill study one-shot-import \
  --manifest jobs/openicpsr_186241/full_one_shot/run-manifest.json \
  --results-dir results \
  --path reports/openicpsr_186241/full_probability_jobs.json --replace
\end{lstlisting}
\newpage
That mapping is passed into the final lifecycle command:

\begin{lstlisting}[style=shell]
zwill study finish \
  --manifest workflow/openicpsr_186241/full/jobs/run-manifest.json \
  --results-dir results/openicpsr_186241/full --replace \
  --audit-out zwill_work/reports/openicpsr_186241/full_audits \
  --out 'zwill_work/reports/openicpsr_186241/full_validation/{survey}' \
  --probability-jobs reports/openicpsr_186241/full_probability_jobs.json \
  --embedder edsl --embedding-model text-embedding-3-small --seed 20260730
\end{lstlisting}
This refuses incomplete results, imports forecasts, renders prompt audits, and
fits the matched XGBoost and one-shot comparisons. Finally, stable study tables
are produced without paper code reopening internal validation files:

\begin{lstlisting}[style=shell]
zwill study summarize \
  --manifest ../config/openicpsr_186241_kalro_study.json \
  --bundle-template 'zwill_work/reports/openicpsr_186241/full_validation/{survey}' \
  --out zwill_work/reports/openicpsr_186241/full_summary
\end{lstlisting}

Each of the \KALROQuestions{} questions is held out in its own approved plan.
For that target, both the twin and XGBoost receive the same 200 sampled
respondents, the same profile fields, and every admissible non-target answer in
the battery (capped at eight for the twelve-item battery). Thus neither method
is trained or prompted with any target answer, and the twin is reconstructed
for every held-out question. The run contains \KALROCalls{} calls including
\KALROQuestions{} same-model one-shot forecasts, costs \KALROCost, has no
malformed prediction rows, and produces zero leakage flags in all three
batteries.

Across questions, GPT-5.5 twins improve on exact-context conditional XGBoost by
\KALROConditionalDelta{} NLL (95\% battery-bootstrap interval
$[\KALROConditionalLo,\KALROConditionalHi]$), with lower NLL on
\KALROConditionalWins{} of 20 questions and in all three battery aggregates.
This comparison is favorable but should be read alongside the stronger
same-model population forecast. Relative to one-shot GPT-5.5, the overall
difference is \KALROOneShotDelta{}
($[\KALROOneShotLo,\KALROOneShotHi]$): the interval crosses zero, so the two
methods are statistically indistinguishable overall. The farming-practices
battery favors twins, while the two four-item input batteries are effectively
tied. Twins also remain worse than the oracle empirical marginal by
\KALROMarginalDelta{} NLL. Their implied marginal distributions have
option-level correlation \KALROMarginalCorrelation{} and MAE
\KALROMarginalMAE. The result therefore strengthens external validity and the
case against a simple benchmark-memorization account, but it is not a universal
win over a strong population-level model.

\begin{figure}[ht]
\centering
\includegraphics[width=0.88\textwidth]{../output/figures/kalro_extension.pdf}
\caption{Leave-one-question-out external validation in three KALRO batteries.
Points are battery macro-average GPT-5.5 twin NLL minus the named benchmark;
bars are paired 95\% bootstrap intervals. Negative differences favor twins.}
\label{fig:kalro-extension}
\end{figure}

\begin{table}[ht]
\centering
\scriptsize
\resizebox{\textwidth}{!}{\begin{tabular}{lrrrlrl}\toprule
Battery & Questions & Flags & $\Delta$ NLL$_{XGB}$ & 95\% CI & $\Delta$ NLL$_{1shot}$ & 95\% CI \\
\midrule
KALRO input use & 4 & 0 & -0.655 & [-0.716, -0.598] & 0.005 & [-0.023, 0.030] \\
KALRO farming practices & 12 & 0 & -0.357 & [-0.383, -0.328] & -0.017 & [-0.031, -0.004] \\
KALRO input awareness & 4 & 0 & -0.029 & [-0.064, 0.009] & 0.007 & [-0.023, 0.033] \\
\bottomrule\end{tabular}
}
\caption{KALRO external-validation results against exact-context XGBoost and
the same-model no-persona one-shot baseline. Negative differences favor
respondent twins.}
\label{tab:kalro-extension}
\end{table}

\clearpage
\subsection{Performance relative to uniform}
Uniform prediction assigns $1/|\mathcal A_j|$ to every option and therefore
uses neither population prevalence nor respondent information. It is a weak
reference, but an informative diagnostic: failure to beat it means that a
method's confidence is actively harmful relative to complete ignorance. We
define question-level NLL skill as
\[
 S^{U}_{j}=1-\frac{\mathrm{NLL}_{j,\mathrm{twin}}}
                       {\log |\mathcal A_j|}.
\]
Positive values favor the twin. In the by-question design, GPT-5.5 has lower
NLL than uniform on \UniformQuestionsBetter\ of 69 targets and positive macro
skill in 14 of 15 batteries. Averaging equally across questions, the NLL
difference is \UniformNLLDelta\ (95\% question-cluster bootstrap interval
$[\UniformNLLLo,\UniformNLLHi]$). The corresponding mean Brier difference is
$-0.162$ (95\% CI $[-0.184,-0.140]$). Thus the result is not driven merely by
comparison with a difficult learned benchmark: the forecasts also dominate a
no-information probability vector on nearly every target.

Uniform is not the decisive benchmark because survey options are rarely
equally prevalent. A model can beat uniform by knowing only that one response
is common. The conditional XGBoost comparison remains the respondent-level
gate; uniform performance establishes basic forecast competence, not digital-
twin validity.

Figure~\ref{fig:benchmark-ladder} puts all three references on the same
question-level scale. Each point is one held-out target and the short black bar
is the median. Moving from uniform to the empirical marginal and then to
conditional XGBoost progressively removes gains attributable to knowing option
counts or population prevalence. The remaining negative distribution against
XGBoost is therefore the relevant evidence for respondent-specific signal; its
wider dispersion also makes clear why survey-level intervals and target-level
diagnostics are both needed.

\begin{figure}[ht]
\centering
\includegraphics[width=0.78\textwidth]{../output/figures/benchmark_ladder.pdf}
\caption{By-question GPT-5.5 NLL minus benchmark NLL for all 69 held-out
targets. Negative values favor the twin. The benchmarks use, respectively, no
information, the observed target distribution, and the same respondents and
prompt-level context as the twin. Horizontal bars mark medians.}
\label{fig:benchmark-ladder}
\end{figure}

\subsection{Semantic proximity and question difficulty}
Across the combined 107-question collection, targets closer on average to their
observed context questions tend to be easier for the twins within the same
battery, although the expanded interval now narrowly includes zero. The coefficient of similarity in the NLL
regression is \SemanticNLLCoef{} (battery-bootstrap 95\% interval
$[\SemanticNLLLo,\SemanticNLLHi]$). On the more interpretable 0.1-similarity
scale, this is approximately a 0.100 reduction in mean NLL. Accuracy moves in
the same direction: its coefficient is 0.670 (95\% interval $[0.065,1.299]$),
or roughly 6.7 percentage points per 0.1 similarity.

The stronger test is whether closeness especially helps the language-model
twin rather than making the prediction problem easier for both methods. Here
the evidence is weak. The coefficient for twin advantage over exact-context
XGBoost is \SemanticAdvCoef{}, with a wide interval
$[\SemanticAdvLo,\SemanticAdvHi]$ that includes zero. The maximum-similarity
measure is also inconclusive for both absolute NLL and relative advantage.
Because XGBoost itself uses the same semantic embedding family, this pattern is
coherent: related context appears informative, but this sample does not show
that the twin extracts disproportionately more value from semantic closeness.

Figure~\ref{fig:semantic-proximity} displays the within-battery variation behind
the estimates. The left panel shows the suggestive ``closer is easier''
claim; the right panel shows why it should not be restated as a comparative
mechanism claim. The LOESS curves do not reveal a stable threshold that would
overturn the linear summary, and their wide bands---especially near the sparse
extremes---make the limited support visible. Outliers have considerable leverage
with three to five targets per battery, and no causal interpretation is
warranted.

\begin{figure}[ht]
\centering
\includegraphics[width=0.96\textwidth]{../output/figures/semantic_proximity.pdf}
\caption{Question-level semantic proximity and predictive performance after
removing battery means. Proximity is the mean cosine similarity between the
held-out question and the context questions actually shown across respondent
prompts. Lines are the corresponding battery-fixed-effect slopes. Higher twin
advantage means lower twin NLL relative to XGBoost. Colored curves are
local-linear LOESS fits with pointwise 95\% whole-battery bootstrap bands.}
\label{fig:semantic-proximity}
\end{figure}

\subsection{Aggregate marginal recovery}
Individual forecasts also imply a marginal distribution. For option $a$ of
question $j$, we average the 200 respondent forecasts,
\[
 \widehat\pi_{ja}=\frac{1}{n_j}\sum_i\widehat p_{ija},
 \qquad
 \pi_{ja}=\frac{1}{n_j}\sum_i\mathbf 1\{y_{ij}=a\}.
\]
Across 295 question--option cells in the 69 by-question targets, the correlation
between $\widehat\pi_{ja}$ and $\pi_{ja}$ is \MarginalCorrelation\ and the mean
absolute option-share error is \MarginalMAE. The modal option agrees on
\MarginalTopAgreement\ of questions; mean question-level Jensen--Shannon
divergence is 0.012. Figure~\ref{fig:marginal-fit} shows that fit directly.

\begin{figure}[ht]
\centering
\includegraphics[width=0.62\textwidth]{../output/figures/marginal_fit.pdf}
\caption{Observed option shares and shares implied by averaging GPT-5.5's 200
by-question respondent forecasts. Each point is one labeled option on one of
69 held-out questions; the dashed line is equality.}
\label{fig:marginal-fit}
\end{figure}

Aggregate recovery and individual prediction answer different questions. The
strong marginal fit shows that averaging the twins yields plausible survey
distributions, but it cannot by itself show that probability mass is assigned
to the correct people. The NLL comparison to exact-context XGBoost supplies
that harder evidence.

Figure~\ref{fig:calibration-comparison} reports a complementary probability
diagnostic. Each line joins the average by-question expected calibration error
of GPT-5.5 and conditional XGBoost within a battery. Calibration is not a
substitute for a proper scoring-rule comparison: a method can be calibrated yet
uninformative. It does reveal whether an NLL advantage is accompanied by
systematically unreliable confidence. The mixed within-battery ordering here
argues against reducing the main result to a universal calibration advantage.

\begin{figure}[ht]
\centering
\includegraphics[width=0.86\textwidth]{../output/figures/calibration_comparison.pdf}
\caption{Mean expected calibration error across by-question runs within each
survey battery. Lower values are better. Lines pair GPT-5.5 twins and the
exact-information conditional XGBoost benchmark evaluated on the same 200
respondents.}
\label{fig:calibration-comparison}
\end{figure}

\subsection{Same-model one-shot marginal benchmark}
The uniform and empirical-marginal references bracket two extremes: no
knowledge and direct access to the realized target distribution. We add a
deployable intermediate comparison by asking GPT-5.5 once per target for the
population response distribution using only the question wording and legal
options. The model receives no respondent profile, no observed answers, and no
empirical target frequencies. If its one-shot distribution is $q_j(a)$, its
score on the same 200 respondents used for by-question validation is
\[
 \mathrm{NLL}_{j,\mathrm{one\mbox{-}shot}}
   =-\frac{1}{n_j}\sum_i \log q_j(y_{ij}),
 \qquad
 \Delta^{OS}_j=\mathrm{NLL}_{j,\mathrm{twin}}
                   -\mathrm{NLL}_{j,\mathrm{one\mbox{-}shot}}.
\]
Negative $\Delta^{OS}_j$ isolates the value of respondent-specific context
while holding the language model and target wording fixed.

The workflow intentionally exports one probability job per battery, rather
than one monolithic cross-study job. Each package contains that battery's
three to five target questions. A newly ingested battery therefore adds one
new export, result, import receipt, and report without rerunning any completed
battery. The Make pattern obtains the target list from the study manifest and
executes the following \zwill--Expected Parrot sequence for each battery:

\begin{lstlisting}[style=shell]
zwill edsl build --survey $survey --target probability-job \
  --questions "$questions" --model openai:gpt-5.5 \
  --path jobs/one-shot/$survey/one_shot_jobs.ep
\end{lstlisting}

\begin{lstlisting}[style=shell]
ep run --jobs jobs/one-shot/$survey/one_shot_jobs.ep --background
\end{lstlisting}

\begin{lstlisting}[style=shell]
zwill prob-results import --survey $survey \
  --input-path results/one-shot/$survey/one_shot_results.ep --replace
\end{lstlisting}

\begin{lstlisting}[style=shell]
zwill prob-results report --survey $survey --job-id $job_id \
  --format json --path reports/one-shot/$survey/marginals.json
\end{lstlisting}

The run added \OneShotCalls\ forecasts across all 15 batteries at an observed
cost of \OneShotCost; all rows imported without issue. Mean question-level
option-share MAE is \OneShotMAE\ for the one-shot distributions and
\OneShotTwinMAE\ for distributions obtained by averaging the 200 corresponding
twins. Figure~\ref{fig:one-shot-marginals} plots both against observed shares.
The twins' aggregate advantage is modest, which is expected: respondent
context is primarily intended to allocate probability to the right people,
not simply to recover population prevalence.

\begin{figure}[ht]
\centering
\includegraphics[width=0.72\textwidth]{../output/figures/one_shot_marginal_fit.pdf}
\caption{Observed option shares versus the GPT-5.5 one-shot distribution and
the average of 200 GPT-5.5 respondent-specific forecasts. Both series cover
the same 295 options from 69 held-out questions.}
\label{fig:one-shot-marginals}
\end{figure}

At the individual level, twins have lower NLL on \OneShotTwinWins\ of 69
questions. The equally weighted mean difference is \OneShotNLLDelta. The
battery-level differences in Figure~\ref{fig:one-shot-lift} favor twins in
\OneShotBatteryWins\ of 15 batteries. Municipal vignettes are the exception:
the no-persona forecast has lower mean NLL there, consistent with that
battery's inconclusive twin--XGBoost result. Thus GPT-5.5's general knowledge
of survey marginals explains some predictive performance, but not the full
benefit usually obtained when the same model receives a respondent's observed
answers and profile.

\begin{figure}[ht]
\centering
\includegraphics[width=0.86\textwidth]{../output/figures/one_shot_nll_lift.pdf}
\caption{Battery-level mean NLL for GPT-5.5 twins minus the same model's
no-persona one-shot distribution, evaluated on the identical by-question
respondent samples. Negative values favor respondent-specific context.}
\label{fig:one-shot-lift}
\end{figure}

\begin{table}[ht]
\centering
\scriptsize
\begin{tabular}{lrrrr}\toprule
Battery & Questions & One-shot MAE & Twin MAE & $\Delta$ NLL \\
\midrule
Brazil digital credit & 5 & 0.061 & 0.029 & -0.216 \\
Climate policy (W158) & 4 & 0.038 & 0.037 & -0.207 \\
Democratic norms & 5 & 0.052 & 0.045 & -0.228 \\
Gender comparisons (W154) & 3 & 0.044 & 0.043 & -0.040 \\
Human vs. AI (W152) & 5 & 0.075 & 0.036 & -0.379 \\
Korea mobility/fairness & 5 & 0.094 & 0.066 & -0.219 \\
Local-news providers (W141) & 5 & 0.023 & 0.049 & -0.180 \\
Municipal vignettes & 4 & 0.074 & 0.094 & 0.195 \\
Origin-group attitudes & 5 & 0.069 & 0.050 & -0.146 \\
Remote work (flagged) & 5 & 0.080 & 0.067 & -0.202 \\
Social-media use (W150) & 5 & 0.022 & 0.036 & -0.061 \\
Spirituality (W132) & 5 & 0.103 & 0.058 & -0.150 \\
Views of social media (W163) & 3 & 0.037 & 0.046 & -0.314 \\
Wellbeing & 5 & 0.056 & 0.031 & -0.288 \\
Worker skills (W157) & 5 & 0.036 & 0.026 & -0.274 \\
\bottomrule\end{tabular}

\caption{One-shot and twin-implied aggregate error, plus individual-level NLL
lift, by independently executable battery. MAE averages equally over questions
within battery; $\Delta$ NLL is twin minus one-shot.}
\label{tab:one-shot}
\end{table}

\subsection{Robustness across language models}
Table~\ref{tab:model-comparison} reports the predeclared five-battery subset.
All three models beat uniform and the empirical-marginal oracle on average,
and all three clear the XGBoost NLL interval in every selected battery.
GPT-5.5 has the lowest mean survey-level NLL in all five batteries and the
largest average skill, although these model-to-model differences were not the
primary paired gate. Gemini 3.5 Flash and Claude Sonnet 4.6 are close to each
other on mean NLL, with neither uniformly dominating at the question level.

\begin{table}[ht]
\centering
\small
\begin{tabular}{lrrrrr}\toprule
Model & Questions & Mean NLL & Skill vs. uniform & Skill vs. marginal & XGB passes \\
\midrule
GPT-5.5 & 22 & 0.935 & 0.268 & 0.137 & 5/5 \\
Claude Sonnet 4.6 & 22 & 0.988 & 0.225 & 0.086 & 5/5 \\
Gemini 3.5 Flash & 22 & 0.986 & 0.232 & 0.094 & 5/5 \\
\bottomrule\end{tabular}

\caption{Matched by-question model comparison over five batteries and 22
targets. Skill is $1-\mathrm{NLL}_{model}/\mathrm{NLL}_{reference}$, averaged
over survey-level macros. ``XGB passes'' counts batteries whose entire paired
95\% interval for model minus XGBoost NLL is below zero.}
\label{tab:model-comparison}
\end{table}

The cross-provider replication matters because digital-twin performance could
otherwise be an idiosyncrasy of one model's instruction following or
probability calibration. The common positive result suggests that the
construction carries transferable signal, while GPT-5.5's consistent edge
also shows that model choice remains consequential.

Figure~\ref{fig:model-comparison} shows the underlying battery-level values
rather than only their cross-battery averages. All four points within a row use
the identical targets and respondents. This view exposes both the common
advantage over XGBoost and the amount of between-battery variation that the
pooled table necessarily compresses.

\begin{figure}[ht]
\centering
\includegraphics[width=0.88\textwidth]{../output/figures/model_comparison.pdf}
\caption{Matched mean NLL for three language-model twins and conditional
XGBoost on the five-battery robustness subset. Lower values are better. Each
battery contains the same 200 respondents and the same by-question information
sets for every method.}
\label{fig:model-comparison}
\end{figure}

\subsection{Audit findings and corrective run}
The implementation keys every baseline feature to the exact observed-answer
list stored for the twin case and fingerprints that information policy in the
baseline job identifier. This makes both respondent membership and available
answers mechanically auditable for every comparison.

The Gallup remote-work audit found a near-duplicate pair: the February 1
work-from-home capacity target and a current-job capacity context item had
Cramér's V of 0.939 among 1,878 respondents. We stopped interpretation of that
target, registered a new construction approach with an explicit target-specific
exclusion, approved and executed a 200-prediction corrective plan, and inspected
the rendered prompts. Those prompts contained seven substantive context
answers and omitted the flagged current-job item. The generic association
audit nevertheless continued to flag the pair because it scans the survey's
eligible question pairs rather than only the rendered prompt. To remain
conservative, we exclude the target---and the remote-work battery's apparent
aggregate pass---from the clean-survey headline. Fourteen primary survey
bundles otherwise had zero leakage flags.

\section{What the evidence supports}
The study supports a positive but bounded predictive conclusion. Conditional
on the selected public batteries, target questions, 200-respondent samples,
model versions, and prompt construction, GPT-5.5 battery-context twins assign
better probabilities than the implemented XGBoost transfer baseline in all 14
clean simultaneous-holdout batteries and 13 of 14 by-question batteries. The
agreement between NLL and Brier gates makes that comparison less dependent on
one scoring rule. It does not expand the population to which the comparison
applies.

This is not evidence that twins dominate every supervised alternative. The
baseline deliberately honors the twin's small information budget and must
transfer across questions without target labels. A production system with
labeled examples for the new target estimates a different and easier task.
Nor do held-out survey answers validate arbitrary counterfactuals. The evidence
supports respondent-level interpolation within related batteries and motivates
replication across models and constructions, not unqualified deployment.

The finite-sample predictive comparison itself is not heuristic: it has a
defined loss, matched cases, a frozen information policy, and uncertainty over
the evaluated respondents. It provides statistically disciplined evidence that
one predictor performs better than another on the declared tasks. In Hullman et
al.'s taxonomy, the heuristic step would be to use that result as proof of
human--twin interchangeability on unobserved tasks \citep{hullman2026human}.
Our controls rule out several mundane reasons a benchmark can look successful,
including prompt-level target leakage, unequal information, unmatched cases,
malformed-row selection, and weak baselines. Eliminating those failures still
does not establish no pretraining leakage, risk generalization to new scenarios,
or preservation of the moments needed for a downstream estimand.

The appropriate deployment implication is therefore prospective and
calibrated. If twins are used to prioritize questions, stress-test an
instrument, or allocate scarce human measurement, the benchmark provides
evidence that their rankings contain signal, while final claims should still be
checked on human outcomes. If twins are used to estimate a mean, subgroup gap,
regression coefficient, or causal effect, the model predictions should be
treated as auxiliary measurements. A jointly labeled human validation sample
can estimate a rectifier, for example through prediction-powered inference,
rather than assuming model errors average away. Better NLL can improve the
precision available from such an approach; it is not itself the identification
argument \citep{ludwig2025llm}.

\section{Limitations}
First, the 15-battery primary study, five-battery GSS extension, and
three-battery KALRO external validation use GPT-5.5; cross-provider robustness is
measured on a predeclared five-battery subset rather than the full corpus. The
study still uses one construction strategy and one fixed seed per design and
does not estimate variation across seeds, temperatures, or prompt pipelines.
Second, 200 respondents provide
useful pilot evidence but limited power for small effects and rare response
options. The per-question complete-case sample may also differ across targets,
although every twin--baseline comparison uses identical cases. Third, all
baselines use one fixed semantic embedding model; results may differ with other
representations or XGBoost hyperparameters. Fourth,
questions within a battery often share wording and topic. This is a realistic
within-instrument prediction task, but easier than predicting an unrelated
future question. Fifth, complete-case filtering changes the evaluated
population, and survey weights were not used as prompt evidence. Sixth,
statistical leakage screening cannot prove the absence of every semantic
shortcut; prompt inspection remains essential.

Seventh, our use of ``leakage'' is narrower than no training leakage. We verify
that held-out answers and target-revealing fields do not enter the prompt, but
cannot determine whether public questionnaire text, reports, or distributions
appeared in proprietary pretraining corpora. Holding a response out at
inference time is not a temporal or corpus holdout. Results therefore describe
performance on the evaluated finite tasks, not unbiased risk over unseen survey
scenarios.

Eighth, target questions and batteries are not probability samples from a
defined universe of behavioral scenarios. Respondent-bootstrap intervals are
conditional on these items, and the cross-question summaries are descriptive
of the 69 primary targets, 38 GSS extension targets, and 20 KALRO targets.
Generalization over questions would require a declared
scenario population and random sampling or an equivalent design. Finally, we
do not evaluate bias for any downstream mean, subgroup contrast, regression,
or treatment effect. Aggregate marginal agreement does not guarantee
conditional moment preservation. Any inferential use of twin-generated
outcomes requires its own human validation sample, estimand-specific error
analysis, and calibration procedure.

\section{Reproducibility}
The paper directory contains declarative ingestion specifications, the survey
cohort configuration, 15 simultaneous and 127 by-question approved plans, 44 additional
model-comparison plans, exported EDSL jobs, raw Results packages,
import records, run-audit pages, leakage files, conditional-baseline
predictions, bootstrap JSON and SVG artifacts, consolidated report bundles,
stable \texttt{zwill study summarize} CSV/JSON exports---including the
independently sourced KALRO extension---with canonical
aggregate statistics and claim gates, package-generated SVG diagnostics,
realized-cost rollups, and presentation scripts that generate Tables~\ref{tab:cohort},
\ref{tab:validation}, \ref{tab:kalro-extension}, and
\ref{tab:model-comparison} plus their figures. The
main cross-survey facts are also stored as CSV and JSON. Rebuilding refuses to
overwrite the completed \texttt{workdir} unless
the operator explicitly passes \texttt{--fresh}.

\appendix
\section{Executed twin prompt template}
The following is the template exposed by the \zwill{} run-audit artifact for the
by-question jobs. Bracketed double braces are populated from the approved EDSL
scenario. The executed jobs recorded no separate system prompt. The actual
answer is retained as hidden scenario metadata for later scoring, but it is not
one of the template variables and is absent from the rendered user prompt.

\begin{lstlisting}[style=prompt]
You are acting as a digital twin for one survey respondent.

Survey name:
{{ survey_name }}

Survey context:
{{ survey_context }}

Respondent id:
{{ respondent_id }}

Non-survey agent construction material:
{{ agent_material_text }}

Supplemental twin material:
{{ twin_material_text }}

Observed answers from this respondent:
{{ observed_answers_text }}

Held-out question name:
{{ heldout_question_name }}

Held-out question text:
{{ heldout_question_text }}

Held-out response options:
{{ heldout_options_text }}

Use the non-survey agent construction material, supplemental twin
material, and observed answers to infer this respondent's probability
distribution over the held-out response options.

Return only valid JSON. Do not include markdown fences, prose, or comments.

The JSON must have exactly this shape:
{
  "probabilities": [0.17, 0.83],
  "notes": "Brief explanation of the respondent-level probability estimates."
}

The probabilities array must contain one number for each held-out option,
in the same order as the options. Each probability must be between 0 and 1.
The probabilities should sum to 1.
\end{lstlisting}

Although the template includes a respondent identifier for joining artifacts,
the identifier is an opaque study-specific key rather than substantive
personality evidence. In the executed examples, non-survey and supplemental
material explicitly state when nothing was supplied. The predictive evidence
is therefore concentrated in readable profile metadata and the structured
\texttt{observed\_answers\_text} block. Every run-audit page preserves rendered
examples and raw responses, permitting direct verification that target answers,
weights, and target marginals were absent.

\section{Reliability and construction questions for future tests}
The executed study deliberately holds construction fixed. The conceptual
framework nevertheless suggests several preregistered extensions. Encoder
ablations can compare answers-only, metadata-only, shuffled-answer, and other-
respondent contexts. Decoder robustness can vary topic distance, wording,
option order, language, and model family. Prompt stability, predictive entropy,
and respondent-level dispersion are candidate reliability signals, but they
should be evaluated on held-out batteries rather than trusted because they are
intuitively appealing. Elaborate persona prose is likewise a hypothesis, not a
free improvement: any compression of the raw answer record may discard the
very local structure that supports prediction. Each extension can be expressed
as a registered \zwill{} approach and approved plan, then compared using the same
leakage audit and exact-information baseline gate.

\bibliographystyle{plainnat}
\bibliography{references}

\end{document}
